% Einheit 4 – Sachaufgaben
\clearpage
\setcounter{aufgabe}{35}
\hypertarget{e4}{}\einheitenkopf{Einheit 4 von 4 · Sachaufgaben}
\zweigzeile{Hier lernst du, zu Zahlenrätseln und Rechtecken eine quadratische Gleichung aufzustellen, zu lösen und die Lösung auszuwählen, die zur Frage passt · neu in diesem Jahr (an der Oberschule je nach Buch Klasse 9 oder 10) · keine P10-Aufgabe · baut auf: Wurzelziehen (Einheit 1), Satz vom Nullprodukt (Einheit 2), p-q-Formel (Einheit 3)}

\begin{aufgabe}{Ich erkenne, welche Lösung zur Frage passt. Die Gleichung ist schon gelöst. Kreuze an, welche Lösung zur Frage passt – oder beide. Rechne nichts.}
\begin{teile}
\teil Breite eines Rechtecks in cm: \ $x_1 = 4$, \ $x_2 = -9$ \hfill \kreuz{$x_1$}\kreuz{$x_2$}\kreuz{beide}
\teil Anzahl der Personen in einer Gruppe: \ $x_1 = 8$, \ $x_2 = -9$ \hfill \kreuz{$x_1$}\kreuz{$x_2$}\kreuz{beide}
\teil eine Zahl, deren Quadrat $49$ ist: \ $x_1 = 7$, \ $x_2 = -7$ \hfill \kreuz{$x_1$}\kreuz{$x_2$}\kreuz{beide}
\teil Flugzeit eines Balls in Sekunden: \ $t_1 = -0{,}5$, \ $t_2 = 2{,}5$ \hfill \kreuz{$t_1$}\kreuz{$t_2$}\kreuz{beide}
\teil kleinere von zwei aufeinanderfolgenden Zahlen: \ $x_1 = 5$, \ $x_2 = -6$ \hfill \kreuz{$x_1$}\kreuz{$x_2$}\kreuz{beide}
\end{teile}
\end{aufgabe}

\begin{aufgabe}{Ich kann ein Zahlenrätsel mit einer quadratischen Gleichung lösen. Gib alle Zahlen an, die passen.}
\beispiel{&\text{Das Quadrat einer Zahl ist } 625. \\ x^2 &= 625 &&\mid \sqrt{\phantom{x}} \\ x_1 &= 25, \quad x_2 = -25 && \text{Die Zahl ist } 25 \text{ oder } {-25}.}
\begin{teile}
\teil Das Quadrat einer Zahl ist $400$. \quad ($x^2 = 400$) \hfill \leerfeld oder \leerfeld
\teil Das Quadrat einer Zahl ist $900$. \quad ($x^2 = 900$) \hfill \leerfeld oder \leerfeld
\teil Das Quadrat einer Zahl ist $0{,}25$. \quad ($x^2 = 0{,}25$) \hfill \leerfeld oder \leerfeld
\teil Das Quadrat einer Zahl ist $2500$. \quad ($x^2 = 2500$) \hfill \leerfeld oder \leerfeld
\end{teile}
\medskip
Stelle die Gleichung selbst auf und löse sie.
\begin{teile}
\teil Addiert man zum Quadrat einer Zahl $19$, erhält man $100$. Welche Zahl kann es sein?
\teil Das Produkt aus einer Zahl und ihrem Nachfolger ist $56$. Gib beide Zahlenpaare an.
\end{teile}
\end{aufgabe}

\begin{aufgabe}{Ich kann eine Sachaufgabe zum Rechteck mit einer quadratischen Gleichung lösen.}
\begin{minipage}[t]{0.64\linewidth}\vspace{0pt}
Beispiel: Ein Rechteck ist $4\,\text{cm}$ länger als breit und hat den Flächeninhalt $45\,\text{cm}^2$.
\rechnung{x\cdot(x + 4) &= 45 \\ x^2 + 4x - 45 &= 0 \\ x_{1,2} &= -2 \pm \sqrt{4 + 45} = -2 \pm 7 \\ x_1 &= 5, \quad x_2 = -9}
Das Rechteck ist $5\,\text{cm}$ breit und $9\,\text{cm}$ lang; $-9$ ist keine Länge.
\end{minipage}\hfill
\begin{minipage}[t]{0.32\linewidth}\vspace{0pt}
\rechteck[punkte=,seiten={x+4,x,,}]{3.6}{2}
\end{minipage}
\medskip
Ein Gemüsebeet ist $5\,\text{m}$ länger als breit und hat den Flächeninhalt $84\,\text{m}^2$.
\begin{teile}
\teil Zeichne eine Skizze, stelle eine Gleichung auf und löse sie.
\teil Welche Lösung passt nicht zur Frage? Begründe.
\teil Schreibe einen Antwortsatz mit Einheit.
\teil Mache die Probe am Beet: Rechne Länge mal Breite.
\teil Ein quadratisches Foto wird vergrößert: Eine Seite wird um $2\,\text{cm}$ länger, die andere um $5\,\text{cm}$. Das neue Foto hat den Flächeninhalt $88\,\text{cm}^2$. Wie lang war eine Seite des alten Fotos?
\end{teile}
\end{aufgabe}

\begin{aufgabe}{Ich finde den Fehler in einer Sachaufgabe. Ein quadratischer Teich ist von einem $1\,\text{m}$ breiten Weg umgeben. Teich und Weg bedecken zusammen $144\,\text{m}^2$. Nina rechnet so:}
\rechnung{(x + 2)^2 &= 144 &&\mid \sqrt{\phantom{x}} \\ x + 2 &= 12 \quad\text{oder}\quad x + 2 = -12 \\ x_1 &= 10, \quad x_2 = -14 && \text{Der Teich ist } 10\,\text{m oder } {-14}\,\text{m lang.}}
\begin{teile}
\teil Finde den Fehler, benenne ihn und korrigiere die Antwort.
\teil Löse selbst: Ein quadratisches Schwimmbecken ist von einem $3\,\text{m}$ breiten Plattenweg umgeben. Becken und Weg bedecken zusammen $400\,\text{m}^2$. Wie lang ist eine Seite des Beckens?
\end{teile}
\end{aufgabe}

\begin{aufgabe}{Ich kann begründen, wann eine negative Lösung zur Frage passt.}
\begin{teile}
\teil Die Gleichung $x^2 = 196$ gehört zu zwei Aufgaben: „Das Quadrat einer Zahl ist $196$.“ und „Ein Quadrat hat den Flächeninhalt $196\,\text{cm}^2$.“ Begründe, warum die erste Aufgabe zwei Antworten hat, die zweite nur eine.
\teil Lukas sagt: „Bei Sachaufgaben ist die negative Lösung immer falsch.“ Stimmt das? Begründe mit einem Beispiel.
\end{teile}
\end{aufgabe}
